%------------------------------------------------------------------------------%
\documentclass[a4paper,12pt,fleqn]{article}

\usepackage{calculo_varias_variaveis-1}

\ShowAnswers

%------------------------------------------------------------------------------%
\begin{document}

\makeHeader{CFVVI}{Prova 1 -- Turma 7}{22/09/2025}

% 01
%------------------------------------------------------------------------------%
\begin{multicols}{2}
\raggedbottom
\exercise{50\%}
  Uma partícula descreve uma trajetória senoidal descrita
  pelas equações paramétricas
  \begin{gather*}
    x(t) = 4 - 3t \\
    y(t) = 6\sin\left(\frac{\pi t}{2}\right) - 2 \qquad\quad
    0 \leq t \leq 2
  \end{gather*}
  \begin{enumerate}[label=\alph*),itemsep=0pt,topsep=0pt]
    \item\!\! [20\%] calcule o vetor tangente à curva
    \item\!\! [10\%] avalie o vetor tangente em \\
                     $t=0$, $t=1$ e em $t=2$
    \item\!\! [20\%] esboce a curva e os vetores \\
                     tangentes calculados
                     (divida os vetores por 3)
  \end{enumerate}
\columnbreak
~

\begin{questiononly}
  \hfill\includegraphics{axis_xy}
\end{questiononly}
\begin{answeronly}
  \hfill\includegraphics{T7_B_parametrica}
\end{answeronly}
\end{multicols}
\clearpagequestiononly

\begin{answer}
\textbf{a)} \quad
O vetor tangente é
\(
  v(t) = \Vetor{\frac{dx}{dt}}{\frac{dy}{dt}}
\)

Avaliando cada derivada

\begin{gather*}
\frac{dx}{dt}
  = \frac{d}{dt}\left[4 - 3t\right]
  = -3
  \\[4mm]
\frac{dy}{dt}
  = \frac{d}{dt}\left[6\sin\left(\frac{\pi t}{2}\right) - 2\right]
  = 6\frac{d}{dt}\sin\left(\frac{\pi t}{2}\right)
  = 6\cos\left(\frac{\pi t}{2}\right)\frac{d}{dt}\left(\frac{\pi t}{2}\right)
  = 6\cos\left(\frac{\pi t}{2}\right)\frac{\pi}{2}
  = 3\pi\cos\left(\frac{\pi t}{2}\right)
\end{gather*}

\medskip
Portanto,
\(
  v(t) = \Vetor{-3}{3\pi\cos\left(\frac{\pi t}{2}\right)}
\)

\clearpage
\textbf{b)} \quad
Avaliando o vetor tangente nos pontos $t=0$, $t=1$ e em $t=2$
\begin{multicols}{3}
\setlength{\mathindent}{0pt}
\begin{align*}
  v(0)
  & = \Vetor{-3}{3\pi\cos\left(\frac{\pi t}{2}\right)} \evalat{0}{} \\[2mm]
  & = \Vetor{-3}{3\pi\cos\left(0\right)} \\[2mm]
  & = \vetor{-3}{3\pi\times 1} \\[2mm]
  & = \vetor{-3}{3\pi}
\end{align*}

\newcolumn
\begin{align*}
  v(1)
  & = \Vetor{-3}{3\pi\cos\left(\frac{\pi t}{2}\right)} \evalat{1}{} \\[2mm]
  & = \Vetor{-3}{3\pi\cos\left(\frac{\pi}{2}\right)} \\[2mm]
  & = \Vetor{-3}{3\pi\times 0} \\[2mm]
  & = \vetor{-3}{0}
\end{align*}

\newcolumn
\begin{align*}
  v(2)
  & = \Vetor{-3}{3\pi\cos\left(\frac{\pi t}{2}\right)} \evalat{2}{} \\[2mm]
  & = \Vetor{-3}{3\pi\cos\left(\pi\right)} \\[2mm]
  & = \Vetor{-3}{3\pi(-1)} \\[2mm]
  & = \vetor{-3}{-3\pi}
\end{align*}
\end{multicols}

\bigskip
\textbf{b)} \quad
Calculando os pontos associados a $t=0$, $t=1$ e em $t=2$
\[
  x(0)
  = \left(4 - 3t\right)\evalat{0}{}
  = 4 - 3\times 0
  = 4
\]
\[
  y(0)
  = \left(6\sin\left(\frac{\pi t}{2}\right) - 2\right)\evalat{0}{}
  = 6\sin\left(0\right) - 2
  = -2
\]

\medskip
\[
  x(1)
  = \left(4 - 3t\right)\evalat{1}{}
  = 4 - 3\times 1
  = 1
\]
\[
  y(1)
  = \left(6\sin\left(\frac{\pi t}{2}\right) - 2\right)\evalat{1}{}
  = 6\sin\left(\frac{\pi}{2}\right) - 2
  = 6 \times 1 - 2
  = 4
\]

\medskip
\[
  x(2)
  = \left(4 - 3t\right)\evalat{2}{}
  = 4 - 3\times 2
  = -2
\]
\[
  y(2)
  = \left(6\sin\left(\frac{\pi t}{2}\right) - 2\right)\evalat{2}{}
  = 6\sin\left(\pi\right) - 2
  = 6 \times 0 - 2
  = -2
\]
\end{answer}

% 02
%------------------------------------------------------------------------------%
\exercise{50\%}
Considerando a função
\(
  f(x, y) = \ln\left(16-2x^2-2y^2\right)
\)
\begin{enumerate}[label=(\alph*),itemsep=0pt]
  \item \ [10\%] determine o domínio de $f$
  \item \ [10\%] esboce o domínio
  \item \ [20\%] determine todas as curvas de nível de $f$
  \item \ [10\%] determine e esboce a curva de nível que passa pelo ponto \((2, 1)\)
\end{enumerate}
\begin{questiononly}
\begin{multicols}{2}
  \;\; Domínio
  \vspace{-0.3\baselineskip}
  \begin{center}
    \includegraphics{axis_xy}
  \end{center}
  \;\; Curva de nível
  \vspace{-0.3\baselineskip}
  \begin{center}
    \includegraphics{axis_xy}
  \end{center}
\end{multicols}
\end{questiononly}
\begin{answeronly}
  \begin{multicols}{2}
    \;\; Domínio
    \vspace{-0.3\baselineskip}
    \begin{center}
      \includegraphics{T7_B_dominio}
    \end{center}
    \;\; Curva de nível
    \vspace{-0.3\baselineskip}
    \begin{center}
      \includegraphics{T7_B_curva_nivel}
    \end{center}
  \end{multicols}
\end{answeronly}

\begin{answer}
\textbf{a)} \quad
  Não podemos calcular logaritmo de um número negativo ou do zero,
  portanto, o domínio de $f$ é o conjunto
  de todos os pontos \((x, y)\in\R^2\) tais que
  \[
    16-2x^2-2y^2 > 0
  \]
  Podemos reescrever essa condição como
  \[
    x^2 + y^2 < 8
  \]
  ou seja, os pontos do disco de raio $\sqrt{8}=2\sqrt{2}$ centrado na origem.
  Assim, o domínio é
  \[
    D_f = \{(x, y) \in \R^2 \st x^2 + y^2 < 8\}
  \]

\bigskip
\textbf{c)} \quad
  Para determinarmos todas as curvas de nível de $f$ precisamos determinar sua
  imagem. Observamos que dentro do domínio temos
  \begin{align*}
    0 < 16-2x^2-2y^2             & \leq 16      \\[2mm]
    \ln\left(16-2x^2-2y^2\right) & \leq \ln(16)
  \end{align*}
  Assim, \(\Im_f = \left\{z\in\R \st z \leq \ln(16)\right\}\).
  As curvas de nível de $f$ consistem dos pontos que satisfazem
  \(
    f(x, y) = c
  \)
  para cada \(c\in\Im_f\).
  \begin{align*}
    f(x, y)                      & = c                 \\[2mm]
    \ln\left(16-2x^2-2y^2\right) & = c                 \\[2mm]
    16-2x^2-2y^2                 & = e^c               \\[2mm]
    - 2x^2 - 2y^2                & = e^c - 16          \\[2mm]
    x^2 + y^2                    & = 8 - \frac{e^c}{2}
  \end{align*}
  As curvas de nível são circunferências centradas na origem com raio
  \(
    r = \sqrt{8 - \frac{e^c}{2}}
  \)

\bigskip
\textbf{c)} \quad
A curva de nível que passa pelo ponto $(2, 1)$ tem o valor
\[
  c
  = f(2, 1)
  = \ln\left(16-2x^2-2y^2\right) \evalat{(2, 1)}{}
  = \ln\left(16-2\times 2^2 - 2\times 1^2\right)
  = \ln\left(16 - 8 - 2\right)
  = \ln\left(6\right)
\]
Correspondendo a uma circunferência de raio
\[
  r
  = \sqrt{8 - \frac{e^c}{2}} \; \evalat{\ln(6)}{}
  = \sqrt{8 - \frac{e^{\ln(6)}}{2}}
  = \sqrt{8 - \frac{6}{2}}
  = \sqrt{8 - 3}
  = \sqrt{5}
  \approx 2,236
\]

\end{answer}

%------------------------------------------------------------------------------%
\end{document}
%------------------------------------------------------------------------------%
